\uuid{GGW8}
\titre{Système linéaire à paramètres/Méthode du Pivot de Gauss}
\auteur{hanani ; mzari}
\organisation{exo7}
\licence{CC BY-NC-SA 4.0}
\datecreate{2020-01-20}
\niveau{L1}
\module{Algèbre}
\chapitre{Déterminant, système linéaire}
\sousChapitre{Système linéaire, rang}
\difficulte{3}
\theme{Systèmes d'équations linéaires, Système linéaire à paramètres, Méthode du Pivot de Gauss}
\exo7qcmid{lille-2.systemes.018}

\contenu{
\begin{question}
\qclassification{200.03}
\qtags{motcle=Système linéaire à paramètres/Méthode du Pivot de Gauss}
On considère le système d'équations, d'inconnue $(x,y,z)\in \Rr^3$ et de paramètres des réels $a$ et $m$ :
$$(\mathtt{S})
\left\{\begin{array}{rcc}
2x+y-z&=&2\\
x-y+z&=&4\\
3x+3y-z&=&4m \\
mx-y+z&=&2a+2. \end{array}\right.$$
Quelles sont les assertions vraies ?
\begin{answers}
\good{$(\mathtt{S})\Leftrightarrow
\left\{\begin{array}{rcc}
x-y+z&=&4\\
y-z&=&-2\\
z&=&2m \\
0&=&m-a. \\
\end{array}\right.$}
\bad{Si $m=1$ et $a=-1$, $(\mathtt{S})$ admet une unique solution.}
\bad{Si $m =a$, $(\mathtt{S})$ admet une infinité de solutions.}
\good{Si $m\neq a$, $(\mathtt{S})$ n'admet pas de solution.}
\end{answers}
\begin{explanations}
$$(\mathtt{S}) \Leftrightarrow
\left\{\begin{array}{rcc}
x-y+z&=&4\\
y-z&=&-2\\
3y-2z&=&2m-6 \\
(m-1)y+(1-m)z&=&2a-4m+2 \\
\end{array}\right. \Leftrightarrow
\left\{\begin{array}{rcc}
x-y+z&=&4\\
y-z&=&-2\\
z&=&2m \\
0&=&m-a. \\
\end{array}\right.$$
Si $m=a$, $(\mathtt{S})$ admet une unique solution et si $m\neq a$, $(\mathtt{S})$ n'admet pas de solution.
\end{explanations}
\end{question}
}
